\begin{claim}{AC.43}{Proved limit of a rank-one overconvergence approach}
Let $U/\mathbf F_q$ be smooth and geometrically connected, and let
$\mathcal L$ be a rank-one unit-root overconvergent $F$-isocrystal
on $U$, with coefficients in a fixed finite extension of
$\mathbf Q_p$. Write $b_x$ for its Frobenius eigenvalue at a closed
point $x$. For every fixed integer $c\ge0$, the finite valuations
\[
 \{v_p(b_x-1):x\in|U|,\ v_p(\deg x)\le c,\ b_x\ne1\}
                                                        \tag{AC.43.1}
\]
are uniformly bounded. This assertion would apply to a character
family only after its return eigenvalues had been identified with
such an overconvergent object.

There is a concrete obstruction to that identification in the
scalar example $C_0:y^2+y=x^5$ over $\mathbf F_{16}$. Its ordinary
character unit roots at the order-three parameters and the order-thirteen
parameter of AC.31 cannot be interpolated by one rank-one coefficient
object with finite geometric monodromy on a common geometrically
connected base over $\mathbf F_{16}$ containing those parameter
points with their indicated closed-point degrees two and three.
In particular no rank-one
unit-root overconvergent $F$-isocrystal can have all those prescribed
return eigenvalues. This conclusion concerns the proposed interpolation;
it does not assert that an intrinsic arithmetic $F$-isocrystal on
the ordinary character locus has been constructed.
\end{claim}
\begin{proof}
By \cite[Proposition~8.6, p.~42]{AC_HartlPal2025}, a tensor power
$\mathcal L^{\otimes N}$ is constant. If its constant Frobenius
eigenvalue is $c_0$, then
\[
                         b_x^N=c_0^{\deg x}.           \tag{AC.43.2}
\]
All these quantities are $p$-adic units. Suppose first that $c_0$
is not a root of unity. In a finite splitting field, let $d$ be the
order of its residue. If $d\nmid m$, then
$v_p(c_0^m-1)=0$. Otherwise write $m=dp^a u$ with $p\nmid u$.
The factorization of $X^u-1$ at
$X=c_0^{dp^a}\equiv1$ gives
\[
 v_p(c_0^m-1)=v_p(c_0^{dp^a}-1).
\]
When $v_p(m)\le c$, the latter runs through a finite set of finite
values. Moreover
\[
 v_p(b_x-1)\le v_p(b_x^N-1)=v_p(c_0^{\deg x}-1),
\]
because the other factor of $b_x^N-1$ is integral. If $c_0$ is a
root of unity, (AC.43.2) puts every $b_x$ in one fixed finite set
of roots of unity, which also proves (AC.43.1). This establishes
the asserted conditional bound.

We now prove the obstruction using two genuine character fibers.
On $J_{C_0}$, Frobenius over $\mathbf F_{16}$ is $[-4]$.
Every nonzero point $a_3\in J_{C_0}[3]$ therefore has degree two.
For the Weierstrass point $P=\infty$, the Raynaud theta curve is
$C_0-P$ and contains no nonzero three-torsion: a function with
divisor $3Q-3P$ would have pole order three at $P$, whereas
$L(3P)=\langle1,x\rangle$. The corresponding pointed cubic
cover over $\mathbf F_{256}$ consequently has ordinary new part,
by the rank-one character criterion. The pointed hyperelliptic
involution commutes with Frobenius and interchanges the two
nontrivial cubic characters. Their common polynomial thus has
coefficients in
$\mathbf Z[\zeta_3]\cap\mathbf R=\mathbf Z$. It is
\[
                         1+s_3T+256T^2,
 \qquad s_3\in\mathbf Z,
\]
and has an ordinary unit root $\beta$ satisfying
$\beta^2+s_3\beta+256=0$. The Weil bound and ordinariness imply
$s_3^2<1024$, so $\mathbf Q(\beta)$ is imaginary quadratic.

Let $a_{13}$ be the explicit order-thirteen parameter of AC.31.
Its degree over $\mathbf F_{16}$ is three. A character unit root
$\alpha$ of its pointed cover over $\mathbf F_{4096}$ satisfies
\[
 P(X)=X^4+139X^3+12473X^2+569344X+16777216=0.       \tag{AC.43.3}
\]
More precisely, put $K_0=\mathbf Q(\sqrt{13})$. Its character
quadratic is
\[
 X^2-tX+4096,
 \qquad t=\frac{-139\pm13\sqrt{13}}2,
\]
with discriminant
\[
 d=t^2-16384=\frac{-22009\mp1807\sqrt{13}}2,
 \qquad
 N_{K_0/\mathbf Q}(d)=110486961=39^2\cdot17\cdot4273.
                                                        \tag{AC.43.4}
\]
Both real conjugates of $d$ are negative, so $d$ is not a square
in the real field $K_0$. Also $t=\alpha+4096/\alpha$ belongs to
$\mathbf Q(\alpha)$ and is irrational. Hence
$[\mathbf Q(\alpha):\mathbf Q]=4$ and $K_0$ is a subfield.
The rational number in (AC.43.4) is neither a rational square nor
thirteen times a rational square: its valuation at seventeen is
one. A rational number is a square in $\mathbf Q(\sqrt{13})$
only in those two square classes. Thus $N(d)$ is not a square
in $K_0$. In particular the nontrivial automorphism of $K_0$
cannot extend to $K_0(\sqrt d)$, since such an extension would
make $\overline d/d$, and hence $d\overline d$, a square in $K_0$.
This proves that $\mathbf Q(\alpha)$ is not Galois. It has no
quadratic subfield other than $K_0$: two different quadratic
subfields would generate a biquadratic degree-four field.

Complex conjugation in this CM field fixes $K_0$ and sends
$\alpha$ to $4096/\alpha$. Under the chosen two-adic embedding,
\[
             v_2(\alpha)=0,
 \qquad v_2(4096/\alpha)=12.
\]
No positive power of $\alpha$ can therefore belong to $K_0$.
Since any proper subfield of a degree-four field has degree one
or two, every positive power of $\alpha$ has degree four.

If one rank-one object with finite geometric monodromy had the
specified eigenvalues at these degree-two and degree-three points,
some tensor power would be constant, giving
\[
                     \beta^N=c_0^2,
 \qquad \alpha^N=c_0^3,
 \qquad \alpha^{2N}=\beta^{3N}.
\]
The right-hand side has degree at most two over $\mathbf Q$,
whereas the left-hand side has degree four. This contradiction
proves the interpolation obstruction. For an overconvergent
rank-one unit-root object, the required finite geometric
monodromy follows from the theorem already cited.

The distinction between a full geometric cohomology object and
its unit-root part is essential. The ordinary Legendre family
has overconvergent full cohomology but a unit-root subobject
which is not overconvergent; see
\cite[Theorem~5.4.3.5, Corollary~5.4.3.10 and
Remark~5.4.3.11, pp.~49--52]{AC_Etesse2010}.
In that example the Frobenius coefficient alone can be
overconvergent while its connection is not. Therefore neither
geometric origin nor a finite algebraic formula for the Cartier
coefficient supplies the missing overconvergence hypothesis.
\end{proof}
